\section{System Architectural Models (Modelli OOA)}

La presente sezione descrive l'architettura logica e comportamentale del sistema \textbf{MyAma} mediante i modelli di \textit{Object Oriented Analysis} (OOA) sviluppati in conformità agli standard UML.

% -----------------------------------------------------------
\subsection{Activity Diagrams}
Gli Activity Diagram modellano il flusso di controllo delle attività nei processi cardine del sistema, evidenziando rami condizionali, parallelismi e punti di decisione.

\subsubsection{Activity Diagram: Prenotazione Ritiro a Domicilio}
\begin{figure}[H]
    \centering
    % \includegraphics[width=0.8\textwidth]{figure/activity_ritiro.png}
    \fbox{\parbox[c][6.5cm][c]{0.8\textwidth}{\centering \textsf{[ Inserire qui l'Activity Diagram per la Prenotazione Ritiro a Domicilio esportato da Visual Paradigm ]}}}
    \caption{Activity Diagram --- Processo di Prenotazione Ritiro a Domicilio}
    \label{fig:act_ritiro}
\end{figure}

\subsubsection{Activity Diagram: Gestione ed Esecuzione Ritiro (Autista)}
\begin{figure}[H]
    \centering
    % \includegraphics[width=0.8\textwidth]{figure/activity_autista.png}
    \fbox{\parbox[c][6.5cm][c]{0.8\textwidth}{\centering \textsf{[ Inserire qui l'Activity Diagram per l'Esecuzione Ritiro da parte dell'Autista esportato da Visual Paradigm ]}}}
    \caption{Activity Diagram --- Flusso Esecuzione Ritiro e Registrazione Esito}
    \label{fig:act_autista}
\end{figure}

\subsubsection{Activity Diagram: Conferimento in Sede (Operatore)}
\begin{figure}[H]
    \centering
    % \includegraphics[width=0.8\textwidth]{figure/activity_sede.png}
    \fbox{\parbox[c][6.5cm][c]{0.8\textwidth}{\centering \textsf{[ Inserire qui l'Activity Diagram per il Conferimento in Sede esportato da Visual Paradigm ]}}}
    \caption{Activity Diagram --- Flusso Accettazione e Verifica al Varco}
    \label{fig:act_sede}
\end{figure}

% -----------------------------------------------------------
\newpage
\subsection{Sequence Diagrams}
I Sequence Diagram descrivono l'interazione dinamica temporale tra gli oggetti del sistema durante la realizzazione degli scenari d'uso, applicando il pattern architetturale \textbf{BCE} (\textit{Boundary}, \textit{Control}, \textit{Entity}).

\subsubsection{Sequence Diagram: UC-C01 Prenota Ritiro a Domicilio}
\begin{figure}[H]
    \centering
    % \includegraphics[width=0.9\textwidth]{figure/seq_prenota_ritiro.png}
    \fbox{\parbox[c][7.5cm][c]{0.85\textwidth}{\centering \textsf{[ Inserire qui il Sequence Diagram per la Prenotazione Ritiro a Domicilio (BCE) esportato da Visual Paradigm ]}}}
    \caption{Sequence Diagram (BCE) --- Prenotazione Ritiro a Domicilio}
    \label{fig:seq_ritiro}
\end{figure}

\subsubsection{Sequence Diagram: UC-A02 Registrazione Esito Ritiro}
\begin{figure}[H]
    \centering
    % \includegraphics[width=0.9\textwidth]{figure/seq_esito_ritiro.png}
    \fbox{\parbox[c][7.5cm][c]{0.85\textwidth}{\centering \textsf{[ Inserire qui il Sequence Diagram per la Registrazione Esito Ritiro Autista esportato da Visual Paradigm ]}}}
    \caption{Sequence Diagram (BCE) --- Registrazione Esito Ritiro}
    \label{fig:seq_esito_ritiro}
\end{figure}

\subsubsection{Sequence Diagram: UC-O01 Convalida Conferimento in Sede}
\begin{figure}[H]
    \centering
    % \includegraphics[width=0.9\textwidth]{figure/seq_conferimento_sede.png}
    \fbox{\parbox[c][7.5cm][c]{0.85\textwidth}{\centering \textsf{[ Inserire qui il Sequence Diagram per la Convalida Conferimento in Sede esportato da Visual Paradigm ]}}}
    \caption{Sequence Diagram (BCE) --- Convalida Conferimento in Sede}
    \label{fig:seq_conferimento_sede}
\end{figure}

% -----------------------------------------------------------
\newpage
\subsection{Class Diagrams}
Il Class Diagram descrive la struttura statica del sistema in termini di classi, attributi, metodi e relazioni strutturali (associazioni, aggregazioni, composizioni ed ereditarietà).

\subsubsection{Class Diagram Unrefined (Modello Preliminare)}
Il modello \textit{Unrefined} riflette le entità di dominio individuate nella fase iniziale di analisi concettuale, focalizzandosi sulle relazioni fondamentali prima dell'introduzione dei pattern e delle ottimizzazioni di visibilità.

\begin{figure}[H]
    \centering
    % \includegraphics[width=0.95\textwidth]{figure/class_unrefined.png}
    \fbox{\parbox[c][8.5cm][c]{0.9\textwidth}{\centering \textsf{[ Inserire qui il Class Diagram Unrefined esportato da Visual Paradigm ]}}}
    \caption{Class Diagram --- Versione Unrefined (Analisi di Dominio)}
    \label{fig:class_unrefined}
\end{figure}

\subsubsection{Class Diagram Refined (Modello Consolidato)}
Il modello \textit{Refined} integra le classi Boundary e Control emerse dai Sequence Diagram, esplicita le firme complete delle operazioni, i tipi di ritorno, i vincoli di molteplicità e la visibilità dei membri.

\begin{figure}[H]
    \centering
    % \includegraphics[width=0.95\textwidth]{figure/class_refined.png}
    \fbox{\parbox[c][8.5cm][c]{0.9\textwidth}{\centering \textsf{[ Inserire qui il Class Diagram Refined esportato da Visual Paradigm ]}}}
    \caption{Class Diagram --- Versione Refined (Specifica Consolidata)}
    \label{fig:class_refined}
\end{figure}
